% =====================================================================
%  Inline-First with Sentinel Promotion
%  A pattern for sparse stores whose cells usually hold one value.
%  arXiv-style preprint. Self-contained: no external .bib required.
%  Build:  pdflatex pattern && pdflatex pattern
% =====================================================================
\documentclass[11pt,a4paper]{article}

\usepackage[T1]{fontenc}
\usepackage[utf8]{inputenc}
\usepackage[margin=1in]{geometry}
\usepackage{amsmath}
\usepackage{amssymb}
\usepackage{amsthm}
\usepackage{booktabs}
\usepackage{xcolor}
\usepackage[hidelinks]{hyperref}

\theoremstyle{plain}
\newtheorem{theorem}{Theorem}
\theoremstyle{definition}
\newtheorem{definition}[theorem]{Definition}
\newtheorem{invariant}{Obligation}
\theoremstyle{remark}
\newtheorem{remark}[theorem]{Remark}

\newcommand{\M}{\mathsf{M}}
\newcommand{\dom}{\operatorname{dom}}

\renewcommand{\topfraction}{0.85}
\renewcommand{\textfraction}{0.12}

\title{\textbf{Inline-First with Sentinel Promotion:\\
A Pattern for Sparse Stores Whose Cells Usually Hold One Value}}

\author{%
Guy Korland\thanks{\texttt{guy.korland@falkordb.com}} \and
Roi Lipman\thanks{\texttt{roi.lipman@falkordb.com}} \and
Dvir Dukhan\thanks{\texttt{dvir.dukhan@falkordb.com}} \and
Avi Avni\thanks{\texttt{avi.avni@falkordb.com}}
\\[0.6em]
FalkorDB\\
\small\texttt{https://github.com/FalkorDB/FalkorDB}}

\date{\today}

\begin{document}
\maketitle

\begin{abstract}
A sparse matrix cell holds one value. Stores built on sparse matrices that need
a cell to hold \emph{several} values --- a multigraph's parallel edges, a
node's labels, a coordinate's posting list --- usually answer with a container
per cell, which pays for the general case in every cell including the
overwhelming majority that hold exactly one thing.

This note extracts, names and generalises the alternative used by FalkorDB's
relationship tensor. \emph{Inline-first with sentinel promotion} keeps the
single value directly in the cell, reserves one bit pattern outside the value
domain as a sentinel, and on the transition to a second value moves
\emph{every} value of that cell --- the formerly inline one included --- into a
single shared overflow structure keyed by a packed coordinate. A cell is then
self-describing: reading it tells you whether you are done or where the rest
lives, with no side table and no speculative second probe.

The mechanism is a paragraph. What is worth writing down is everything else:
the two preconditions that decide whether the pattern applies at all; the six
obligations an adopter inherits, of which one is subtle enough to have been the
source of the design's hardest bug and another turned out to be an obligation
one should \emph{decline}; the concrete complexity cost, which is a
transaction-local state space of ten states per cell where a Boolean store has
four, three of them existing only because cells carry values; and the
measurement traps, since two of our own published attributions of cost were
wrong in ways that a plausible story had concealed. We give the pattern's
machine-checked invariant set, a second worked instantiation as a design
sketch, and an explicit account of when not to use it.
\end{abstract}

\noindent\textbf{Keywords:} sparse matrices, hypersparse, sentinel encoding,
data-structure patterns, multi-version concurrency control, delta encoding,
mechanised verification

% =====================================================================
\section{The problem the pattern solves}

A sparse matrix associates a value with each stored coordinate. That is one
value. A great many stores want a coordinate to carry a small \emph{multiset}:
a property multigraph connects an ordered node pair by arbitrarily many edges,
each with its own identity; a node carries labels; an inverted index maps a
term--document pair to positions.

The canonical answer is a container per cell --- the cell holds a pointer,
usually tagged, to a vector of values. It is correct, it is uniform, and it
prices the rare case into the common one. If the multiplicity distribution is
dominated by $1$, almost every cell pays for an allocation, an indirection and
an object header in order to hold a single value that would have fitted in the
cell.

The pattern here is the obvious alternative made safe. Keep the one value in
the cell. Reserve a bit pattern the value domain cannot produce. When a cell
needs a second value, write the sentinel and move the whole multiset --- not
just the new arrival --- to one shared overflow structure. Cells that never
overflow never learn that the overflow exists.

What makes this worth a note rather than a footnote is that the obvious
alternative is only obvious until it has to be transactional, iterable and
countable at the same time. The obligations in \S\ref{sec:obligations} are the
content.

\paragraph{Provenance.} Everything here is extracted from a shipped
implementation --- the per-relationship-type edge store of FalkorDB's Rust
engine --- and its machine-checked model. That instantiation is described and
measured in a companion paper; this note deliberately says as little about
edges as it can, and states which claims are measured, which are modelled, and
which are conjecture.

% =====================================================================
\section{Preconditions}
\label{sec:pre}

The pattern needs exactly two things. Both are cheap to check before
committing to it, and neither is negotiable.

\begin{definition}[Sentinel headroom]
The value domain is strictly smaller than the cell's representation, so at
least one bit pattern is not a legitimate value.
\end{definition}

This is a statement about the \emph{domain}, not about current data. ``No
identifier is ever $2^{64}-1$ in practice'' is not headroom; ``identifiers are
allocated from a counter bounded by the index limit $2^{60}-1$'' is. The
distinction matters because the sentinel's safety is the one property with no
graceful degradation: a value colliding with the sentinel does not perform
badly, it silently reads as a multi-value cell and returns the wrong answer.
An adopter should be able to name the bound and where it is enforced.

\begin{definition}[Multiplicity dominated by one]
The distribution of values per cell has its mass at $1$.
\end{definition}

The pattern's entire advantage is that the common case is free; the pattern's
entire cost is a promotion transition and a second lookup for cells past the
boundary. If a large fraction of cells hold two or more values, a container per
cell or an always-materialised auxiliary structure is the better design and the
extra invariants buy nothing. We know of no shortcut here: this is an empirical
question about the workload, and it is the assumption the whole construction
rests on.

\begin{remark}[The precondition nobody checks]
Multiplicity is usually assumed rather than measured, because the data to
measure it lives in production and the design decision is taken before
production exists. It is worth saying plainly that this is the pattern's
largest risk and that it is a risk about \emph{deployments}, not about data
structures.
\end{remark}

% =====================================================================
\section{The construction}
\label{sec:construction}

Three components, of which the third is empty in the case the pattern is for.

\begin{description}
  \item[The primary store.] A sparse matrix whose cells are wide enough for a
    value. A cell holding a non-sentinel value \emph{is} the answer.
  \item[The sentinel.] One reserved value $\M$ outside the domain. A cell
    holding $\M$ announces that this coordinate has more than one value and
    that none of them is here.
  \item[The overflow.] One shared structure holding the values of every
    promoted coordinate, keyed by a packed encoding $\kappa$ of the
    coordinate. In the instantiation it is a second hypersparse matrix whose
    row index is $\kappa(\text{src},\text{dst})$ and whose columns are the
    values, so a coordinate's multiset is one row and arrives in sorted order
    for free.
\end{description}

\paragraph{Promotion.} A coordinate at one value that gains a second writes
$\M$ into the cell and moves \emph{both} values to the overflow. Moving the
incumbent is the step that is easy to skip and must not be: it is what makes
the cell self-describing, and \S\ref{sec:obligations} states it as an
obligation rather than an implementation detail.

\paragraph{Demotion.} A coordinate falling back to one value restores that
value inline and empties its overflow row. Whether demotion should be eager or
hysteretic is a policy question; we measured it for the instantiation and found
the upside of hysteresis bounded at about a quarter of an oscillation's cost
against a permanent space charge, and did not build it.

\paragraph{The packed key is a design decision with teeth.} $\kappa$ must be
injective over the coordinates the store admits. Packing two $32$-bit halves
into a $64$-bit row index is the obvious choice and it silently imposes a
$2^{32}$ ceiling on each coordinate axis. That is a real limit, not a
theoretical one, for any store that churns identifiers. Adopters should choose
$\kappa$ with the ceiling in mind and treat widening it as a serialisation
change, because it is one.

% =====================================================================
\section{The obligations you inherit}
\label{sec:obligations}

This is the part that does not transfer by reading the mechanism. Write
$\varepsilon(p)$ for the effective value at coordinate $p$, $\mathrm{Vals}(p)$
for its multiset, and let the store be transactional with a committed base $m$,
a pending-addition layer $dp$ and a pending-deletion mask $dm$.

\begin{invariant}[Mask domain]\label{ob:mask}
$dm \subseteq \dom(m)$. The deletion mask only ever marks committed cells.
\end{invariant}

\begin{invariant}[Purity]\label{ob:pure}
$\dom(dp) \cap dm = \emptyset$. A cell is never simultaneously pending and
deleted.
\end{invariant}

\begin{invariant}[Cancel-to-clean]\label{ob:cancel}
$\forall p \in \dom(dp) \cap \dom(m): dp(p) \neq m(p)$. An operation that would
restore a coordinate's committed value drops its deltas instead of recording a
delta that happens to agree.
\end{invariant}

\begin{invariant}[Promotion completeness]\label{ob:promo}
$\varepsilon(p) = \M \iff |\mathrm{Vals}(p)| \ge 2$, and then \emph{every}
value of $p$ is in the overflow at $\kappa(p)$. Conversely, if
$|\mathrm{Vals}(p)| \le 1$, the overflow at $\kappa(p)$ is empty.
\end{invariant}

\begin{invariant}[Count agreement]\label{ob:count}
Any cached count of promoted coordinates equals
$|\{p : \varepsilon(p) = \M\}|$.
\end{invariant}

\begin{invariant}[Mirror agreement]\label{ob:mirror}
Any secondary index over the same coordinates has exactly the effective
pattern as its domain.
\end{invariant}

\subsection{Which of these actually bite}

Obligations~\ref{ob:mask} and~\ref{ob:pure} are the standard delta hygiene any
transactional layer needs and the pattern does not make them harder. The other
four are where the pattern's character shows.

\paragraph{Promotion completeness is the whole value proposition.}
It is what makes a cell self-describing. A non-sentinel value is \emph{the}
value, never one of several, so no reader ever consults the overflow
speculatively --- and that is precisely the cost that an always-materialised
auxiliary structure pays on every read forever. If an implementation leaves the
incumbent inline on promotion, ``read the cell'' stops being a complete answer
and the pattern degenerates into a container-per-cell design with worse
ergonomics.

It also has a dual worth stating: the \emph{converse} clause, that a
non-promoted coordinate's overflow row is empty, is what lets the number of
promoted coordinates be read off the overflow's structure rather than tracked.
That is the observation that retires Obligation~\ref{ob:count} below.

\paragraph{Cancel-to-clean is the one that is easy to get wrong.}
It is the least obvious obligation and the most load-bearing, and it is
invisible until something counts. Without it, a transaction that promotes a
coordinate, demotes it, and promotes it again can leave $dp(p) = \M$ shadowing
$m(p) = \M$. Reads are unaffected --- both layers say the same thing --- so
tests that read pass. What breaks is cardinality: the number of
\emph{syntactically} shadowed coordinates diverges from the number of
\emph{semantically} shadowed ones, and any closed-form count over the layers
goes wrong. The rule is that the representation must be canonical: a delta
exists only where the state genuinely differs from the committed state.

Two consequences follow. Counting stays exact and closed-form, and the
per-coordinate state space stays bounded --- without canonicalisation there is
no finite automaton to reason about, because ``pending, and equal to
committed'' multiplies every state.

\paragraph{Count agreement is an obligation to decline.}
This one changed status during the work, and the change is the most
transferable lesson in the note. It began as a genuine invariant over a
maintained counter: every mutation path had to remember to move it. It was the
only piece of state resting on a discipline rather than a structural argument,
and it was consequently the only piece with a class of bug --- a path that
forgets --- that no local reasoning rules out.

Promotion completeness makes the counter unnecessary. A promoted coordinate is
exactly one whose overflow row is non-empty, and a structure that stores only
non-empty rows can report how many it has as metadata. Deriving the quantity
instead of maintaining it deleted the invariant, deleted the discipline, and
deleted the per-mutation obligations attached to it; in the mechanisation the
clause became a definitional identity holding by \texttt{rfl}.

The general form of the lesson: \emph{in a design whose invariants make a
quantity structurally determined, caching that quantity converts a theorem into
an obligation.} Look for the derived form before accepting a counter. We state
the obligation here anyway, because it is what an adopter would need if they
ever cache the count again --- and because an adopter who does not know it
exists will cache it.

\paragraph{Mirror agreement is a reminder that the sentinel is not a value.}
Any secondary structure --- a transpose, a reverse index --- must mirror the
\emph{effective pattern}, and a promoted cell is one coordinate, not several.
Mirrors that carry structure only (rather than duplicating values) pay one
primary lookup per hit to recover whether the coordinate is inline or promoted.
That is a real cost and a deliberate trade: it buys not storing every value
twice.

% =====================================================================
\section{The cost you take on}
\label{sec:cost}

Patterns that only list benefits are advertisements. The concrete cost here is
a strictly larger transaction-local state space, and it is not small.

A Boolean transactional cell has four states: absent, present, pending-add,
pending-delete. A value-carrying cell under this pattern has ten, and
\emph{three of them exist only because cells carry values} --- they are the
states in which a pending layer \emph{shadows} a live committed value rather
than adding to or subtracting from a pattern. In-place update is the operation
that forces them: a Boolean layer never needs to say ``this coordinate is still
present, but its value is now different''.

The whole cancel-to-clean normalisation exists for the same reason. So the
honest accounting is: three of ten states, plus a canonicalisation rule that
every mutation path must respect, is the price of carrying values in cells. A
design that keeps the primary store Boolean and puts all values in an
always-materialised auxiliary structure needs none of it --- and pays a second
lookup on every read, forever, including the single-value reads that are the
overwhelming majority.

That is the trade in one sentence, and it is the trade an adopter should make
consciously.

% =====================================================================
\section{What it buys, and how to find out}
\label{sec:measured}

We report the instantiation's numbers as evidence that the trade can come out
favourably, not as numbers that transfer. Against a pointer-tagged
container-per-cell design in the same system, measured at the data-structure
boundary:

\begin{itemize}
  \item At multiplicity $1$ there is \emph{nothing to choose} between the two
    on space: both cost $23.1$ bytes per coordinate. The pattern's advantage is
    not that the common case is cheaper than a container design's common case
    in bytes --- it is that the common case involves no allocation, no
    indirection, and no second structure.
  \item Promoting a coordinate costs $46.1$ bytes against the container
    design's $272.6$, a factor of $5.9$. The entire space difference is a fixed
    per-coordinate charge, not a per-value one: past the second value, the
    container design is marginally \emph{cheaper} per value ($8.04$ bytes
    against $9.09$).
  \item Reads of promoted coordinates are dearer, not cheaper. Scanning a
    hypersparse overflow row costs more per value than walking a container.
\end{itemize}

The shape to take away is that the pattern wins on the $\mu \approx 1$ mass and
the promotion transition, and loses on high-multiplicity access. That is
exactly what the preconditions predict, and it is why \S\ref{sec:pre} is not
boilerplate.

\subsection{How the measurement goes wrong}
\label{sec:measurement}

Three traps, each of which caught us, and each of which we would expect to
catch an adopter evaluating this pattern against a container design.

\paragraph{Differencing whole-system measurements can invert an ordering.}
We estimated the per-value cost of both designs by differencing engine-level
measurements at different multiplicities, and concluded the pattern's overflow
row was cheaper per value than a container walk --- with a crossover above
which it wins outright. Measured at the data-structure boundary the two slopes
are $124$ and $33$, the other way round. Both engine-level figures were high,
as one would expect from differencing, but they were high by $3.6\times$ and
$26\times$ --- and errors that do not share a common factor do not cancel in a
ratio. There is no crossover. \emph{A ratio of two differences of whole-system
measurements can come out backwards}, and no amount of repetition detects it.

\paragraph{A plausible mechanism will survive if nothing asks it to account for
the total.} We attributed a $2\times$ deficit on one operation to the pattern
lacking a specialised entry point that the container design has, and published
it. The story fit: it explained the direction, and the deficit grew with
multiplicity exactly as the story predicted. Decomposing it before implementing
the fix showed the attributed mechanism accounted for $81$ instructions out of
$1{,}990$. Four fifths was per-call iterator construction --- an implementation
artifact shared with every other read path, and the fix for it was worth far
more than the fix we had proposed. The check that would have caught this
immediately is arithmetic: \emph{do the named costs sum to the measured total?}

\paragraph{Instruction counts and time disagree, in both directions.}
Instruction counts are reproducible and immune to a noisy host, which is why we
used them. They are also blind to residency, and the blindness is not a
conservative bias. Sweeping the working set from far inside cache to far
outside it, the pattern's multi-value read costs a flat $3.8$--$4.0\times$ the
container design's in instructions at every size --- while in time it costs
$4.67\times$ in cache and $1.88\times$ out of it, because the container
design's larger footprint leaves cache first. At multiplicity $1$ the error
runs the other way. An evaluation on instruction counts alone will misjudge
this pattern against a container design at realistic sizes, and it will do so
in a direction that depends on the multiplicity.

That last one has a constructive corollary worth stating: \emph{the pattern's
space advantage is also a latency advantage}, and the space constants predict
where. Ours put the container design past a $100$\,MB working set at about
$3.4\times10^5$ coordinates and the pattern's at about $1.5\times10^6$; the
measured points where each fell out of cache bracketed both predictions.

% =====================================================================
\section{A second instantiation, as a sketch}
\label{sec:labels}

To show the pattern is not edge-shaped, here is a candidate we have
\emph{not} built, stated at the level of detail its status deserves.

A property-graph engine stores which labels a node carries. A natural
representation is a Boolean matrix $\text{node} \times \text{label}$: a node
with one label has one entry. Multiplicity is dominated by $1$ --- most nodes
carry one label --- so precondition~2 holds, and label identifiers are drawn
from a small counter, so precondition~1 holds with enormous headroom.

The inline-first form replaces the matrix with a vector
$\text{node} \to \text{label}$, sentinel-promoted into a shared overflow for
multi-labelled nodes. ``Which labels does this node carry'' becomes a point
read rather than a row scan, and ``is this node labelled $L$'' becomes a
comparison.

Three things would have to be established before believing it, and they are the
same three for any candidate:

\begin{enumerate}
  \item \textbf{Does it actually save space?} Probably not, and this is the
    interesting part. A Boolean matrix entry stores an index and no value at
    all; the inline form stores an index and a $64$-bit value. The pattern's
    space case rests on displacing a \emph{container}, and there is no
    container here to displace. A candidate whose baseline is already
    allocation-free may be a lookup win and a space loss.
  \item \textbf{Which direction dominates the workload?} Label queries are
    largely the reverse direction --- all nodes with label $L$ --- which the
    inline form answers badly and which is typically served by a separate
    per-label structure anyway. The pattern should be applied to the direction
    whose multiplicity is dominated by $1$, and only there.
  \item \textbf{Is the promotion transition on a hot path?} The pattern moves
    cost to the boundary crossing. A workload that adds and removes labels
    frequently pays it repeatedly.
\end{enumerate}

We record this as a sketch precisely because the honest answer to (1) is
unfavourable-looking, and a pattern note that only exhibits instantiations that
work is not much of a guide.

% =====================================================================
\section{Mechanisation transfers better than code}
\label{sec:mech}

The instantiation's obligations, its per-coordinate automaton and its two main
theorems are machine-checked in Lean 4 against a model transcribed from the
implementation --- $7{,}071$ lines and $460$ theorems, no \texttt{sorry} and no
bespoke axiom. Two things about that are relevant to an adopter rather than to
us.

First, \emph{what the proofs found}. Two branches of the shipped code were
unreachable, one of which would have fabricated a valid identifier had the
invariant beneath it ever broken --- a fallback worse than a crash, because the
value it produced was indistinguishable from a real one. Two fast paths were
proved to agree with their slow paths. None of this came from proving the
headline theorems; it came from being forced to state the obligations precisely
enough to discharge them.

Second, \emph{what it does not establish}. The proofs bind a model, not the
binary. The correspondence is a transcription that a human checked. An adopter
inheriting the obligation set inherits a set of statements that are known to be
consistent and sufficient --- which is the useful thing, and is not the same as
their implementation being correct.

The transferable artifact is therefore the obligation set of
\S\ref{sec:obligations}, stated in a form that can be re-mechanised against a
different model, rather than the proof scripts.

% =====================================================================
\section{When not to use it}
\label{sec:dont}

\begin{itemize}
  \item \textbf{Multiplicity is not dominated by one.} Use a container per cell
    or an always-materialised auxiliary. The obligations buy nothing and the
    ten-state algebra is pure cost.
  \item \textbf{The value domain fills the cell.} No headroom, no sentinel. Do
    not manufacture headroom by declaring a value ``impossible'' unless it is
    impossible by construction and you can name where.
  \item \textbf{The baseline has no container to displace.} If cells already
    store a bare index with no value array, as in a Boolean pattern matrix, the
    pattern is a lookup optimisation and probably a space regression
    (\S\ref{sec:labels}).
  \item \textbf{High-multiplicity access dominates.} Reads of promoted
    coordinates are the pattern's weak side, and measurably so.
  \item \textbf{The store is not transactional and never will be.} Then most of
    \S\ref{sec:obligations} is vacuous, the pattern reduces to the mechanism,
    and it needs no note --- which is a good outcome, just not one requiring
    this one.
\end{itemize}

% =====================================================================
\section{Conclusion}

Inline-first with sentinel promotion is a small idea with a large obligation
set. The mechanism --- keep one value in the cell, reserve a sentinel, move the
whole multiset on promotion --- takes a paragraph and is arguably folklore. What
does not take a paragraph is: the two preconditions that decide whether it
applies; promotion completeness, which is what makes a cell self-describing and
is the entire reason readers never probe speculatively; cancel-to-clean, which
is invisible until something counts and is where the design's hardest bug came
from; the discovery that one of the six obligations should be declined rather
than discharged, because the invariants already determine the quantity; the
ten-state cost, three states of which are the price of values in cells; and a
set of measurement traps that caught the authors twice in print.

We offer the pattern on those terms. An adopter who takes the mechanism without
the obligations will get a store that reads correctly and counts wrongly, and
will find out late.

\end{document}
